\documentclass[a4paper,12pt]{article}
\usepackage[T2A]{fontenc}
\usepackage[utf8]{inputenc}
\usepackage[russian]{babel}
\usepackage{amsmath,amssymb,mathtools}
\usepackage{helvet}
\renewcommand{\familydefault}{\sfdefault}
\usepackage{icomma}
\usepackage{geometry}
\geometry{left=24mm,right=24mm,top=23mm,bottom=24mm}
\usepackage{microtype}
\usepackage{xcolor}
\usepackage{tikz}
\usetikzlibrary{calc,arrows.meta,positioning,shapes.geometric}
\usepackage{booktabs,tabularx}
\usepackage{enumitem}
\usepackage{hyperref}
\usepackage{fancyhdr}
\usepackage{setspace}
\usepackage{needspace}
\definecolor{accent}{HTML}{147A7C}
\definecolor{darkaccent}{HTML}{155457}
\definecolor{textgray}{HTML}{30383A}
\definecolor{muted}{HTML}{687576}
\definecolor{soft}{HTML}{D5E7E5}
\hypersetup{colorlinks=true,linkcolor=darkaccent,urlcolor=darkaccent,
  pdfauthor={Артём Александрович Лёвин},
  pdftitle={Чек-лист. Многочлены и метод интервалов. 14.10.26}}
\setstretch{1.22}
\setlength{\parindent}{0pt}
\setlength{\parskip}{5pt}
\setlist[itemize]{leftmargin=6mm,itemsep=2pt,topsep=3pt}
\setlist[enumerate]{leftmargin=8mm,itemsep=6pt,topsep=3pt}
\setlength{\headheight}{15pt}
\pagestyle{fancy}
\fancyhf{}
\fancyhead[L]{\small\color{muted}Алгебра \,$\cdot$\, многочлены и неравенства}
\fancyhead[R]{\small\color{muted}14.10.2026}
\fancyfoot[C]{\small\color{muted}\thepage}
\renewcommand{\headrulewidth}{0pt}
\newcommand{\topic}[1]{\Needspace{3\baselineskip}\vspace{8pt}{\Large\bfseries\color{darkaccent}#1}\par\vspace{2pt}\noindent\color{soft}\rule{\linewidth}{0.6pt}\color{textgray}\par\vspace{4pt}}
\newcommand{\subtopic}[1]{\Needspace{2\baselineskip}\vspace{5pt}{\large\bfseries\color{darkaccent}#1}\par}
\newcommand{\idea}[1]{\noindent\makebox[6mm][l]{\color{accent}$\square$}\begin{minipage}[t]{\dimexpr\linewidth-7mm}\raggedright #1\end{minipage}\par\vspace{3pt}}
\newcommand{\important}[1]{\par\textbf{Важно.} #1\par}
\newcommand{\exercise}[2]{\item \textbf{#1}\par #2\par}
\AtBeginDocument{\color{textgray}}

\begin{document}
\thispagestyle{empty}
\begin{center}
\vspace*{26mm}
{\small\color{accent}\bfseries АЛГЕБРА \,$\cdot$\, СИСТЕМАТИЗАЦИЯ МЕТОДОВ}\par
\vspace{13mm}
{\fontsize{31}{37}\selectfont\bfseries\color{darkaccent}Чек-лист}\par
\vspace{8mm}
{\fontsize{21}{29}\selectfont\bfseries Многочлены}\par
\vspace{2mm}
{\fontsize{17}{25}\selectfont и метод интервалов}\par
\vspace{13mm}
{\large Деление с остатком и теорема Безу}\par
{\large Схема Горнера и разложение на множители}\par
{\large Кратность корней, рациональные неравенства}\par
{\large Замена переменной в уравнениях}\par
\vspace{16mm}
{\color{accent}\rule{38mm}{0.7pt}}\par
\vspace{9mm}
{\large\bfseries Ксения Клыкова}\par
\vspace{3mm}
{\normalsize 14 октября 2026 года}\par
\vfill
{\small\color{muted}Лёвин Артём Александрович}\par
{\small\color{darkaccent}\bfseries эксклюзивно для Ксении Клыковой}\par
\vspace{19mm}
\end{center}
\begin{tikzpicture}[remember picture,overlay]
  \draw[accent!55,line width=1.1pt]
    ([xshift=25mm,yshift=25mm]current page.south west)
    .. controls ([xshift=66mm,yshift=41mm]current page.south west)
    and ([xshift=-65mm,yshift=11mm]current page.south east)
    .. ([xshift=-24mm,yshift=27mm]current page.south east);
\end{tikzpicture}
\clearpage

\topic{1. Деление многочленов с остатком}
\idea{Для многочленов $A(x)$ и $B(x)\ne0$ существуют единственные многочлены $Q(x)$ и $R(x)$ такие, что
\[ A(x)=B(x)Q(x)+R(x),\qquad R(x)=0\ \text{или}\ \deg R<\deg B. \]
Здесь $Q$~--- частное, $R$~--- остаток.}
\idea{При делении уголком очередной член частного выбирайте по \textbf{старшим членам} текущего остатка и делителя. Умножайте, вычитайте, повторяйте, пока степень остатка не станет меньше степени делителя.}

\textbf{Пример из занятия.} Разделим
\[
A(x)=2x^3-5x^2+5x-2
\quad\text{на}\quad B(x)=x^2-x-1.
\]
Первый член частного $2x$:
\[
\begin{aligned}
A(x)-2xB(x)&=(2x^3-5x^2+5x-2)\\
 &\quad -(2x^3-2x^2-2x)\\
&=-3x^2+7x-2.
\end{aligned}
\]
Следующий член частного $-3$:
\[
(-3x^2+7x-2)-(-3)B(x)=4x-5.
\]
Итак, $Q(x)=2x-3$, $R(x)=4x-5$, причём $\deg R=1<2$.
\[
\boxed{A(x)=(x^2-x-1)(2x-3)+(4x-5).}
\]

\topic{2. Теорема Безу: остаток без деления уголком}
\idea{При делении $P(x)$ на $x-a$ остаток равен \textbf{$P(a)$}. Следовательно, $x-a$ является множителем $P(x)$ тогда и только тогда, когда $P(a)=0$.}
\textbf{Пример.} Найдём остаток при делении
$P(x)=3x^3-4x^2-2x+1$ на $x-1$:
\[
P(1)=3-4-2+1=\boxed{-2}.
\]
\important{В теореме Безу делитель должен иметь вид $x-a$. Для делителя $2x-3$ его корень $a=\tfrac32$, а остаток по-прежнему равен $P(\tfrac32)$.}

\topic{3. Разложение на множители и схема Горнера}
\idea{Если $P(x)$ имеет целые коэффициенты и ненулевой свободный член, то \textbf{любой целый корень} делит свободный член. При произвольном целом старшем коэффициенте это условие остаётся необходимым для целого корня.}
\idea{Подбирайте кандидат $a$, вычисляйте $P(a)$. Если результат ноль, делите $P(x)$ на $x-a$ уголком или по схеме Горнера, после чего продолжайте разложение многочлена меньшей степени.}

\textbf{Пример из занятия.} Для $P(x)=-3x^2+15x+72$ проверим $a=-3$:
\[
P(-3)=-3\cdot9-45+72=0.
\]
Схема Горнера для коэффициентов $(-3;15;72)$:
\[
\renewcommand{\arraystretch}{1.4}
\begin{array}{r|rrr}
-3 &-3&15&72\\
   &  &+9&-72\\ \hline
   &-3&24&\mathbf0
\end{array}
\]
\[
\boxed{-3x^2+15x+72=(x+3)(-3x+24)=-3(x+3)(x-8).}
\]
\important{Последнее число в схеме~--- остаток, предыдущие числа~--- коэффициенты частного по убыванию степеней. При нулевом остатке найденный множитель подтверждён.}

\topic{4. Метод интервалов и кратности корней}
\idea{Перенесите всё в одну сторону и, по возможности, представьте левую часть как произведение множителей или одну дробь. Найдите действительные нули числителя и знаменателя.}
\idea{При переходе через точку с \textbf{нечётной} кратностью знак меняется; с \textbf{чётной} кратностью знак сохраняется. Если сокращаете общие множители, первоначальные запреты $Q(x)\ne0$ остаются.}
\idea{В строгом неравенстве нули числителя исключаются. В нестрогом они включаются только тогда, когда знаменатель в точке отличен от нуля. \textbf{Нули знаменателя не включаются никогда.}}

\textbf{Простой пример.} Решите $(x+3)(x+1)(x-5)>0$. Точки $-3,-1,5$ простые. При $x>5$ знак положительный; при переходе через каждую точку знак меняется. Поэтому
\[
\boxed{x\in(-3;-1)\cup(5;+\infty).}
\]

\subtopic{Чётная кратность и запрещённые точки}
\textbf{Пример из занятия.} Решите
\[
\frac{(x-1)^{10}(x+2)^9}{2x-3}\ge0.
\]
Критические точки: $-2$ (нечётная кратность $9$), $1$ (чётная кратность $10$), $\frac32$ (знаменатель, нечётная кратность $1$). В точке $x=0$ выражение отрицательно.
\[
\renewcommand{\arraystretch}{1.5}
\begin{array}{c|cccc}
\text{Интервал}&(-\infty;-2)&(-2;1)&(1;\frac32)&(\frac32;+\infty)\\ \hline
\text{Знак}&+&-&-&+
\end{array}
\]
Точка $-2$ включается как нуль числителя; точка $1$ включается отдельно, хотя знак при переходе не меняется. Точка $\frac32$ исключается:
\[
\boxed{x\in(-\infty;-2]\cup\{1\}\cup\left(\frac32;+\infty\right).}
\]
\important{Если перенести множитель $(x-1)^{10}$ в знаменатель, то $x=1$ станет запрещённой точкой. Даже чётная степень в знаменателе создаёт разрыв области определения.}

\subtopic{Как распутывать большие степени}
\textbf{Пример.} Для выражения
\[
\frac{(x^2-2x+1)^{20}(x^3-x^2+3x-3)}{(4x-10)^{19}(x-10)^{24}}
\]
перед расстановкой знаков разложите:
\[
\begin{aligned}
(x^2-2x+1)^{20}&=(x-1)^{40},\\
x^3-x^2+3x-3&=(x-1)(x^2+3),
\end{aligned}
\]
поэтому весь числитель равен $(x-1)^{41}(x^2+3)$. Для вещественных $x$ множитель $x^2+3$ всегда положителен. Критические точки: $1$ (нечётная кратность $41$), $\frac52$ (нечётная кратность $19$) и $10$ (чётная кратность $24$, но точка запрещена). При знаке $<0$ ответ:
\[
\boxed{x\in\left(1;\frac52\right).}
\]


\topic{5. Замена переменной в многочленных уравнениях}
\idea{Если в уравнении повторяется одно выражение $f(x)$, положите $t=f(x)$, решите более простое уравнение для $t$, а затем выполните \textbf{обратную замену} для каждого найденного значения.}
\textbf{Пример из занятия.} Решите
\[
(x^3-x^2)^2+8(x^3-x^2)-48=0.
\]
Положим $t=x^3-x^2$. Тогда
\[
t^2+8t-48=(t+12)(t-4)=0,
\quad t=-12\ \text{или}\ t=4.
\]
После обратной замены получаем два уравнения:
\[
\begin{aligned}
x^3-x^2+12&=0,\\
x^3-x^2-4&=0.
\end{aligned}
\]
Разложим их с помощью найденных целых корней:
\[
\begin{aligned}
x^3-x^2+12&=(x+2)(x^2-3x+6),\\
x^3-x^2-4&=(x-2)(x^2+x+2).
\end{aligned}
\]
Дискриминанты обоих квадратных трёхчленов отрицательны:
$D_1=9-24=-15$, $D_2=1-8=-7$. Поэтому
\[
\boxed{x=-2\ \text{или}\ x=2.}
\]
\idea{Если у всех членов есть общий множитель, вначале вынесите его за скобки. Например, $x^3-3x^2+2x=x(x-1)(x-2)$. Это помогает сразу увидеть часть корней.}

\topic{6. Самопроверка перед записью ответа}
\idea{Проверьте тождество $A=BQ+R$ и условие $\deg R<\deg B$ при делении с остатком.}
\idea{При использовании теоремы Безу пересчитайте $P(a)$; при схеме Горнера проверьте последний столбец.}
\idea{В рациональном неравенстве сначала выпишите все запреты знаменателя, затем проверьте чётность кратностей и строгость знака.}
\idea{После замены переменной решите \emph{все} уравнения обратной замены и проверьте полученные действительные корни.}

\clearpage
\thispagestyle{fancy}
\begin{center}
{\LARGE\bfseries\color{darkaccent}Алгоритм: метод интервалов}\par
\vspace{3mm}
{\small\color{muted}От разложения на множители к ответу}\par
\end{center}
\vspace{8mm}
\begin{center}
\begin{tikzpicture}[
  >=Stealth,
  every node/.style={align=center,font=\small},
  terminal/.style={ellipse,draw=accent,line width=.9pt,fill=soft!25,
    text width=94mm,minimum height=10mm,inner sep=3pt},
  operation/.style={rectangle,rounded corners=2mm,draw=accent!85,line width=.85pt,
    text width=118mm,minimum height=12mm,inner sep=4pt},
  arrow/.style={-Stealth,draw=darkaccent,line width=.85pt}
]
\node[terminal] (a) {Дано неравенство с многочленами};
\node[operation,below=8mm of a] (b) {Перенесите всё в одну сторону;\\разложите числитель и знаменатель на множители};
\node[operation,below=8mm of b] (c) {Найдите действительные корни и их кратности;\\отдельно выпишите запрещённые значения};
\node[operation,below=8mm of c] (d) {Нанесите критические точки на числовую прямую\\в порядке возрастания};
\node[operation,below=8mm of d] (e) {Определите знак на одном интервале\\подстановкой пробного значения};
\node[operation,below=8mm of e] (f) {При переходе через нечётную кратность\\меняйте знак; через чётную~--- сохраняйте};
\node[operation,below=8mm of f] (g) {Выберите интервалы нужного знака;\\добавьте допустимые нули при $\ge$ или $\le$};
\node[terminal,below=8mm of g] (h) {Исключите все нули знаменателя. Запишите ответ};
\foreach \u/\v in {a/b,b/c,c/d,d/e,e/f,f/g,g/h}
  \draw[arrow] (\u.south)--(\v.north);
\end{tikzpicture}
\end{center}
\clearpage

\topic{Тренировочный блок \texorpdfstring{\,$\cdot$\,}{ · } Путь А}
\textit{10 упражнений.} Выполняйте задания последовательно. В задачах на неравенства указывайте полное множество решений. Для каждого ответа полезно сделать короткую проверку.
\begin{enumerate}[label=\textbf{\arabic*.},leftmargin=9mm,itemsep=5pt]
\exercise{Теорема Безу.}{Найдите остаток от деления $P(x)=2x^3-3x^2+4x-5$ на $x-1$.}
\exercise{Деление уголком.}{Разделите $2x^3-x^2-2x-6$ на $x^2-x-2$. Укажите частное и остаток.}
\exercise{Схема Горнера.}{Разложите на множители $2x^2-5x-3$, предварительно проверив, является ли $3$ корнем.}
\exercise{Целые корни.}{Полностью разложите на линейные множители $x^3-4x^2+x+6$.}
\exercise{Простой метод интервалов.}{Решите $(x+2)(x-1)(x-4)>0$.}
\exercise{Нуль чётной кратности.}{Решите \(\displaystyle\frac{(x-1)^2(x+3)}{x-2}\ge0\).}
\exercise{Запрещённая точка чётной кратности.}{Решите \(\displaystyle\frac{(x+1)^3(x-2)^2}{(x-4)(x-1)^2}\le0\).}
\exercise{Большие степени.}{Решите
\[\frac{(x^2+4x+4)^{12}(x^3+2x^2+5x+10)}{(2x-6)^{11}(x-1)^8}<0.\]}
\exercise{Замена переменной.}{Решите уравнение
\[ (x^3+x^2)^2+2(x^3+x^2)-8=0. \]}
\exercise{Итоговая задача.}{Решите неравенство, используя разложение на множители и метод интервалов:
\[(x^3-x^2)^2+8(x^3-x^2)-48\ge0.\]}
\end{enumerate}
\clearpage
\topic{Ответы к тренировочному блоку}
\renewcommand{\arraystretch}{1.65}
\begin{tabularx}{\linewidth}{@{}>{\bfseries}p{12mm}X@{}}
\toprule
№ & Ответ \\ \midrule
1 & $-2$.\\
2 & Частное: $2x+1$; остаток: $3x-4$.\\
3 & $(x-3)(2x+1)$.\\
4 & $(x+1)(x-2)(x-3)$.\\
5 & $(-2;1)\cup(4;+\infty)$.\\
6 & $(-\infty;-3]\cup\{1\}\cup(2;+\infty)$.\\
7 & $[-1;1)\cup(1;4)$.\\
8 & $(-2;1)\cup(1;3)$.\\
9 & $\{-2;1\}$.\\
10 & $(-\infty;-2]\cup[2;+\infty)$.\\
\bottomrule
\end{tabularx}
\vspace{8mm}
\textbf{Контрольные ориентиры.} В заданиях 6--8 отдельно проследите, какие точки являются нулями числителя, а какие запрещены знаменателем. В заданиях 9--10 используйте замену переменной и разложение кубических многочленов.
\end{document}